% 【本文件写什么】api-v0 契约层，叶子，只写语义不写实现
% - crate 路径：os/components/wateros-klog/klog-api/api-v0/src/
% 【事实来源】
% - 上述 crate 下全部 .rs 源文件
% - docs/exports/public-api/ 中对应符号表，以源码为准
% 【不写什么】
% - impl 内部数据结构、汇编、具体算法步骤

\paragraph{KlogRecordMeta}

\code{repr(C)}固定布局记录头：\code{seq}，提交后单调序号、\code{ts\_nsec}、\code{text\_len}、\code{facility}、\code{flags}、\code{level}、\code{caller\_id}。正文另存环内，不裸\code{copy\_to\_user}导出元数据结构。\code{new}构造未提交记录，\code{seq} 由环实现填入；\code{traditional\_level\_char}供读路径格式化为\code{'0'..'7'}。

facility/level常量对齐syslog：\code{LOG\_KERN}、\code{LOG\_USER}；\code{LOG\_EMERG}至\code{LOG\_DEBUG}。

\paragraph{KlogFlags}

位域：\code{CONT}，续行、\code{TRUNC}，截断、\code{USER}，用户态WRITE。\code{empty}/\code{with}/\code{contains}操作\code{u8}编码。

\paragraph{KlogStore trait}

环存储契约，由\code{impl-ringbuf}实现：
\begin{itemize}
 \item \code{append(meta, text) -> AppendResult}：写入记录，实现分配\code{seq}；
 \item \code{stats() -> KlogStats}：提交/丢弃计数与序号边界；
 \item \code{unread\_bytes()}：读游标后未读正文字节总和，\code{SIZE\_UNREAD} 近似；
 \item \code{buffer\_bytes()}：text环容量，即 \code{SIZE\_BUFFER}；
 \item \code{peek\_next\_unread()}：取下一条未读；无未读返回\code{KlogError::NoUnread}；
 \item \code{advance\_read\_cursor(after\_seq)}：消费后推进游标；
 \item \code{clear\_read\_cursor()}：\code{CLEAR}，mark全部已读，物理记录保留。
\end{itemize}

\paragraph{syslog action}

常量\code{SYSLOG\_ACTION\_CLOSE}至\code{SYSLOG\_ACTION\_SIZE\_BUFFER}，0--10与Linux \code{klogctl}对齐。\code{decode\_action}识别已知action；\code{is\_write\_priority}判断WRITE路径，非 0--10 的 \code{type}，priority 编码 level/facility。

\paragraph{错误与结果}

\code{KlogError::NoUnread}，READ返回0、\code{BufferTooSmall}。\code{AppendResult \{ seq, truncated \}}。\code{KlogRecordView}为提交后只读视图，含 \code{meta} 与正文切片。

\paragraph{上下文}

trait方法假定在持锁或等价互斥的上下文中调用；并发语义由impl定义。元数据不直接暴露给用户态指针，导出经\code{format\_traditional}格式化。
